\documentclass[reqno]{amsart}
\usepackage[T1]{fontenc}
\usepackage{amsmath,amssymb,amsthm}
\usepackage{xr-hyper}
\usepackage[hidelinks]{hyperref}
\externaldocument[ec-]{20261124T223000--elliptic-curves-for-the-proof__workbook}
\externaldocument[cp-]{20261125T006500--covolumes-and-periods__workbook}

% Chapter number; \ref{chap} gives it.
\newcounter{chap}
% Chapter numbers, generated from 20261125T012000--grh-workbook__guide.org
% by scripts/make-chapters.py.  Do not edit: change the order of the
% index and rerun the script (./preview-site.sh and CI do this).
%
% Usage in a sheet, after \newcounter{chap}:
%   % Chapter numbers, generated from 20261125T012000--grh-workbook__guide.org
% by scripts/make-chapters.py.  Do not edit: change the order of the
% index and rerun the script (./preview-site.sh and CI do this).
%
% Usage in a sheet, after \newcounter{chap}:
%   % Chapter numbers, generated from 20261125T012000--grh-workbook__guide.org
% by scripts/make-chapters.py.  Do not edit: change the order of the
% index and rerun the script (./preview-site.sh and CI do this).
%
% Usage in a sheet, after \newcounter{chap}:
%   \input{chapters}\chapnumber{<sheet ID>}
\makeatletter
\providecommand\chapters@set[2]{\@namedef{chapters@#1}{#2}}
\chapters@set{20261125T001000}{1}% zeros-to-ramp-sums
\chapters@set{20261125T001100}{2}% inert-primes-and-estimate-b
\chapters@set{20261125T005000}{3}% galois-fourier-and-averaging
\chapters@set{20261125T005500}{4}% columns-from-a-sphere
\chapters@set{20261125T000000}{5}% slope-method
\chapters@set{20261126T030100}{6}% a-twist-for-each-prime
\chapters@set{20261125T003000}{7}% the-chain-and-the-square-root
\chapters@set{20261125T010000}{8}% rank-lemma-and-zero-estimate
\chapters@set{20261125T010500}{9}% poincare-series-and-near-orthogonality
\chapters@set{20261124T223000}{10}% elliptic-curves-for-the-proof
\chapters@set{20261125T009500}{11}% fock-model-and-theta-intertwiners
\chapters@set{20261124T223100}{12}% complex-multiplication-and-inert-reduction
\chapters@set{20261125T008000}{13}% one-cm-determinant
\chapters@set{20261124T214500}{14}% supersingular-theta-divisibility
\chapters@set{20261124T214600}{15}% the-isotropic-plane
\chapters@set{20261125T013700}{16}% pairings-on-a-formal-group-of-height-two
\chapters@set{20261125T013500}{17}% hecke-nondegeneracy-on-the-tree
\chapters@set{20261126T030200}{18}% the-bound-is-attained
\chapters@set{20261125T009000}{19}% kronecker-limit-and-chowla-selberg
\chapters@set{20261125T002000}{20}% definite-theta-lifts
\chapters@set{20261125T015500}{21}% siegel-weil-in-the-smallest-case
\chapters@set{20261125T002100}{22}% definite-theta-lifts-and-norms
\chapters@set{20261125T016500}{23}% hecke-operators-and-theta-lifts
\chapters@set{20261125T016700}{24}% divisibility-along-the-chain
\chapters@set{20261125T004000}{25}% denominators-and-sha
\chapters@set{20261125T004100}{26}% tate-twists-and-sha
\chapters@set{20261126T020430}{27}% the-minus-part-by-gauss-sums
\chapters@set{20261125T014100}{28}% congruence-modules
\chapters@set{20261125T014500}{29}% eisenstein-congruences-and-selmer-classes
\chapters@set{20261125T014700}{30}% pseudocharacters-and-the-length-bound
\chapters@set{20261125T011000}{31}% galois-representations-black-box
\chapters@set{20261126T030300}{32}% the-data
\chapters@set{20261126T020210}{33}% bessel-poincare-series-and-their-gram-operators
\chapters@set{20261126T030400}{34}% the-seesaw-identity
\chapters@set{20261126T030500}{35}% divisibility-of-the-composite-lift
\chapters@set{20261126T030600}{36}% theorem-a-short-forms
\chapters@set{20261126T030700}{37}% the-determinant-of-a-chain
\chapters@set{20261126T030800}{38}% the-cost-of-one-lift
\chapters@set{20261126T030900}{39}% denominators-of-projections-to-lifts
\chapters@set{20261126T031000}{40}% twists-with-few-congruences
\chapters@set{20261126T031100}{41}% theorem-c-and-the-main-inequality
\providecommand\chapnumber[1]{%
  \@ifundefined{chapters@#1}%
    {\PackageError{chapters}{Sheet #1 is not listed in the index}%
      {List it in 20261125T012000--grh-workbook__guide.org\MessageBreak
       and run scripts/make-chapters.py.}}%
    {\edef\thechap{\@nameuse{chapters@#1}}}}
\makeatother
\chapnumber{<sheet ID>}
\makeatletter
\providecommand\chapters@set[2]{\@namedef{chapters@#1}{#2}}
\chapters@set{20261125T001000}{1}% zeros-to-ramp-sums
\chapters@set{20261125T001100}{2}% inert-primes-and-estimate-b
\chapters@set{20261125T005000}{3}% galois-fourier-and-averaging
\chapters@set{20261125T005500}{4}% columns-from-a-sphere
\chapters@set{20261125T000000}{5}% slope-method
\chapters@set{20261126T030100}{6}% a-twist-for-each-prime
\chapters@set{20261125T003000}{7}% the-chain-and-the-square-root
\chapters@set{20261125T010000}{8}% rank-lemma-and-zero-estimate
\chapters@set{20261125T010500}{9}% poincare-series-and-near-orthogonality
\chapters@set{20261124T223000}{10}% elliptic-curves-for-the-proof
\chapters@set{20261125T009500}{11}% fock-model-and-theta-intertwiners
\chapters@set{20261124T223100}{12}% complex-multiplication-and-inert-reduction
\chapters@set{20261125T008000}{13}% one-cm-determinant
\chapters@set{20261124T214500}{14}% supersingular-theta-divisibility
\chapters@set{20261124T214600}{15}% the-isotropic-plane
\chapters@set{20261125T013700}{16}% pairings-on-a-formal-group-of-height-two
\chapters@set{20261125T013500}{17}% hecke-nondegeneracy-on-the-tree
\chapters@set{20261126T030200}{18}% the-bound-is-attained
\chapters@set{20261125T009000}{19}% kronecker-limit-and-chowla-selberg
\chapters@set{20261125T002000}{20}% definite-theta-lifts
\chapters@set{20261125T015500}{21}% siegel-weil-in-the-smallest-case
\chapters@set{20261125T002100}{22}% definite-theta-lifts-and-norms
\chapters@set{20261125T016500}{23}% hecke-operators-and-theta-lifts
\chapters@set{20261125T016700}{24}% divisibility-along-the-chain
\chapters@set{20261125T004000}{25}% denominators-and-sha
\chapters@set{20261125T004100}{26}% tate-twists-and-sha
\chapters@set{20261126T020430}{27}% the-minus-part-by-gauss-sums
\chapters@set{20261125T014100}{28}% congruence-modules
\chapters@set{20261125T014500}{29}% eisenstein-congruences-and-selmer-classes
\chapters@set{20261125T014700}{30}% pseudocharacters-and-the-length-bound
\chapters@set{20261125T011000}{31}% galois-representations-black-box
\chapters@set{20261126T030300}{32}% the-data
\chapters@set{20261126T020210}{33}% bessel-poincare-series-and-their-gram-operators
\chapters@set{20261126T030400}{34}% the-seesaw-identity
\chapters@set{20261126T030500}{35}% divisibility-of-the-composite-lift
\chapters@set{20261126T030600}{36}% theorem-a-short-forms
\chapters@set{20261126T030700}{37}% the-determinant-of-a-chain
\chapters@set{20261126T030800}{38}% the-cost-of-one-lift
\chapters@set{20261126T030900}{39}% denominators-of-projections-to-lifts
\chapters@set{20261126T031000}{40}% twists-with-few-congruences
\chapters@set{20261126T031100}{41}% theorem-c-and-the-main-inequality
\providecommand\chapnumber[1]{%
  \@ifundefined{chapters@#1}%
    {\PackageError{chapters}{Sheet #1 is not listed in the index}%
      {List it in 20261125T012000--grh-workbook__guide.org\MessageBreak
       and run scripts/make-chapters.py.}}%
    {\edef\thechap{\@nameuse{chapters@#1}}}}
\makeatother
\chapnumber{<sheet ID>}
\makeatletter
\providecommand\chapters@set[2]{\@namedef{chapters@#1}{#2}}
\chapters@set{20261125T001000}{1}% zeros-to-ramp-sums
\chapters@set{20261125T001100}{2}% inert-primes-and-estimate-b
\chapters@set{20261125T005000}{3}% galois-fourier-and-averaging
\chapters@set{20261125T005500}{4}% columns-from-a-sphere
\chapters@set{20261125T000000}{5}% slope-method
\chapters@set{20261126T030100}{6}% a-twist-for-each-prime
\chapters@set{20261125T003000}{7}% the-chain-and-the-square-root
\chapters@set{20261125T010000}{8}% rank-lemma-and-zero-estimate
\chapters@set{20261125T010500}{9}% poincare-series-and-near-orthogonality
\chapters@set{20261124T223000}{10}% elliptic-curves-for-the-proof
\chapters@set{20261125T009500}{11}% fock-model-and-theta-intertwiners
\chapters@set{20261124T223100}{12}% complex-multiplication-and-inert-reduction
\chapters@set{20261125T008000}{13}% one-cm-determinant
\chapters@set{20261124T214500}{14}% supersingular-theta-divisibility
\chapters@set{20261124T214600}{15}% the-isotropic-plane
\chapters@set{20261125T013700}{16}% pairings-on-a-formal-group-of-height-two
\chapters@set{20261125T013500}{17}% hecke-nondegeneracy-on-the-tree
\chapters@set{20261126T030200}{18}% the-bound-is-attained
\chapters@set{20261125T009000}{19}% kronecker-limit-and-chowla-selberg
\chapters@set{20261125T002000}{20}% definite-theta-lifts
\chapters@set{20261125T015500}{21}% siegel-weil-in-the-smallest-case
\chapters@set{20261125T002100}{22}% definite-theta-lifts-and-norms
\chapters@set{20261125T016500}{23}% hecke-operators-and-theta-lifts
\chapters@set{20261125T016700}{24}% divisibility-along-the-chain
\chapters@set{20261125T004000}{25}% denominators-and-sha
\chapters@set{20261125T004100}{26}% tate-twists-and-sha
\chapters@set{20261126T020430}{27}% the-minus-part-by-gauss-sums
\chapters@set{20261125T014100}{28}% congruence-modules
\chapters@set{20261125T014500}{29}% eisenstein-congruences-and-selmer-classes
\chapters@set{20261125T014700}{30}% pseudocharacters-and-the-length-bound
\chapters@set{20261125T011000}{31}% galois-representations-black-box
\chapters@set{20261126T030300}{32}% the-data
\chapters@set{20261126T020210}{33}% bessel-poincare-series-and-their-gram-operators
\chapters@set{20261126T030400}{34}% the-seesaw-identity
\chapters@set{20261126T030500}{35}% divisibility-of-the-composite-lift
\chapters@set{20261126T030600}{36}% theorem-a-short-forms
\chapters@set{20261126T030700}{37}% the-determinant-of-a-chain
\chapters@set{20261126T030800}{38}% the-cost-of-one-lift
\chapters@set{20261126T030900}{39}% denominators-of-projections-to-lifts
\chapters@set{20261126T031000}{40}% twists-with-few-congruences
\chapters@set{20261126T031100}{41}% theorem-c-and-the-main-inequality
\providecommand\chapnumber[1]{%
  \@ifundefined{chapters@#1}%
    {\PackageError{chapters}{Sheet #1 is not listed in the index}%
      {List it in 20261125T012000--grh-workbook__guide.org\MessageBreak
       and run scripts/make-chapters.py.}}%
    {\edef\thechap{\@nameuse{chapters@#1}}}}
\makeatother
\chapnumber{20261125T009000}
\renewcommand{\thesection}{\thechap.\arabic{section}}

\theoremstyle{definition}
\newtheorem{problem}{Problem}[section]
\theoremstyle{plain}
\newtheorem{theorem}{Theorem}[section]

\newcommand{\Z}{\mathbb{Z}}
\newcommand{\Q}{\mathbb{Q}}
\newcommand{\R}{\mathbb{R}}
\newcommand{\C}{\mathbb{C}}
\newcommand{\ii}{\mathrm{i}}
\newcommand{\Oc}{\mathcal{O}}
\newcommand{\Beta}{\mathrm{B}}

\setlength{\emergencystretch}{2em}

\title{The Kronecker limit formula and Chowla--Selberg}
\author{}
\date{November 25, 2026}

\begin{document}
\maketitle
\refstepcounter{chap}\label{chap}

\paragraph{Prerequisites.}
Euler's Gamma function: the reflection and duplication formulas, the Beta integral, and \(\psi := \Gamma'/\Gamma\) with \(\psi(1) = -\gamma\).  Poisson summation.  Dirichlet \(L\)-functions: the functional equation \cite[Ch.~9]{Dav}, the class number formula and binary quadratic forms \cite[Ch.~6]{Dav}.  The Hurwitz zeta function \cite[Ch.~12]{Ap}.  Compact Riemann surfaces and their holomorphic differentials, for Section~\ref{sec:periods} only.  From Chapter~\ref{ec-chap}, \emph{Complex tori, theta functions and periods}: Section~\ref{ec-sec:periods}, in particular Problems~\ref{ec-pr:varpi}, \ref{ec-pr:epstein} and~\ref{ec-pr:Omega} and Theorem~\ref{ec-thm:LL}, with Problems~\ref{ec-pr:LL-eta} and~\ref{ec-pr:eta-i}.

Chapter~\ref{ec-chap} took Kronecker's limit formula from the literature and computed \(L'(1,\chi)/L(1,\chi)\) for the character \(\chi\) modulo \(4\) through the discriminant of the lemniscatic lattice.  Here we prove the limit formula and reach the same value through \(s = 0\), where no elliptic functions appear.  Comparing the two routes is the method of Chowla and Selberg \cite{CS}.  It shows that, for the character \(\chi\) of every imaginary quadratic field, \(e^{\gamma - L'(1,\chi)/L(1,\chi)}\) is, up to an algebraic factor, the square of a period of an elliptic curve with complex multiplication divided by \(\pi^2\).

\section{Eisenstein series and Kronecker's limit formula}\label{sec:klf}

What is the constant term at \(s = 1\) of \(\sum'_{\mu \in \Z[\ii]} \lvert\mu\rvert^{-2s}\)?  The answer comes from letting the lattice vary.  For \(\tau = x + \ii y\) with \(y > 0\) and \(\operatorname{Re} s > 1\), let
\[
  E(\tau, s) := \sum_{(m,n) \ne (0,0)} \frac{y^s}{\lvert m + n\tau\rvert^{2s}} .
\]
Let \(q := e^{2\pi\ii\tau}\), let \(\eta(\tau) := q^{1/24}\prod_{n \ge 1}(1 - q^n)\) be Dedekind's function, and let \(\sigma_t(n) := \sum_{e \mid n} e^t\).  For \(z > 0\), let \(K_\nu(z) := \tfrac12 \int_0^\infty t^{\nu - 1} e^{-z(t + t^{-1})/2} \, d t\).

\begin{problem}\label{pr:E-inv}
\(E(g\tau, s) = E(\tau, s)\) for every \(g \in \mathrm{SL}_2(\Z)\).
\end{problem}

\begin{problem}\label{pr:poisson}
For real \(x\), \(y > 0\) and \(\operatorname{Re} s > \tfrac12\),
\[
  \sum_{m \in \Z} \frac{1}{((m + x)^2 + y^2)^{s}} = \frac{\sqrt\pi\, \Gamma(s - \frac12)}{\Gamma(s)\, y^{2s-1}} + \frac{4\pi^s}{\Gamma(s)} \sum_{k \ge 1} \Bigl(\frac{k}{y}\Bigr)^{s - \frac12} \! K_{s - \frac12}(2\pi k y) \cos(2\pi k x) .
\]
\end{problem}

\begin{problem}\label{pr:fourier}
\(E(\tau, s)\) equals
\begin{multline*}
  2\zeta(2s) y^s + \frac{2\sqrt\pi\, \Gamma(s - \frac12)}{\Gamma(s)} \zeta(2s - 1) y^{1-s} \\
  + \frac{8\pi^s \sqrt y}{\Gamma(s)} \sum_{n \ge 1} n^{s - \frac12} \sigma_{1 - 2s}(n) K_{s - \frac12}(2\pi n y) \cos(2\pi n x) .
\end{multline*}
In particular \(E(\tau, s)\) continues meromorphically to \(\operatorname{Re} s > \tfrac12\), with only a simple pole at \(s = 1\).
\end{problem}

\begin{problem}\label{pr:K-half}
\(K_{1/2}(z) = \sqrt{\pi/2z}\; e^{-z}\).
\end{problem}

\begin{problem}\label{pr:q-terms}
At \(s = 1\) the last term in Problem~\ref{pr:fourier} is \(-4\pi \log \bigl\lvert \prod_{n \ge 1} (1 - q^n) \bigr\rvert\).
\end{problem}

\begin{problem}\label{pr:const}
As \(s \to 1\),
\[
  \frac{2\sqrt\pi\, \Gamma(s - \frac12)}{\Gamma(s)} \zeta(2s - 1) y^{1-s} = \frac{\pi}{s - 1} + 2\pi\bigl(\gamma - \log 2 - \tfrac12 \log y\bigr) + O(s - 1) .
\]
\end{problem}

Since \(2\zeta(2) y = \pi^2 y/3 = -4\pi \log\lvert q^{1/24}\rvert\), Problems~\ref{pr:fourier}--\ref{pr:const} prove the following \cite{We}, \cite[Ch.~20, 21]{La}.

\begin{theorem}[Kronecker's first limit formula]\label{thm:klf}
\(E(\tau, s) = \dfrac{\pi}{s - 1} + 2\pi\bigl(\gamma - \log 2 - \log(\sqrt y\, \lvert\eta(\tau)\rvert^2)\bigr) + O(s - 1)\).
\end{theorem}

\begin{problem}\label{pr:eta-mod}
\(\lvert\eta(-1/\tau)\rvert = \lvert\tau\rvert^{1/2} \lvert\eta(\tau)\rvert\).
\end{problem}

\begin{problem}\label{pr:Zi}
Let \(Z(s) := \sum'_{\mu \in \Z[\ii]} \lvert\mu\rvert^{-2s}\), and let \(\chi\) be the non-trivial character modulo \(4\).  Then \(Z(s) = \pi/(s-1) + 2\pi(\gamma - \log 2 - 2\log\eta(\ii)) + O(s-1)\), and
\[
  \frac{L'(1, \chi)}{L(1, \chi)} = \gamma - 2\log 2 - 4\log\eta(\ii) .
\]
\end{problem}

This is Problem~\ref{ec-pr:LL-eta}, now with a proof.  It expresses \(L'/L(1,\chi)\) through a value of a modular form.  To see Gamma values instead, we move to \(s = 0\), where \(L(s,\chi)\) is a finite combination of Hurwitz zeta functions.

\section{Lerch's formula and the functional equation}\label{sec:lerch}

For \(x > 0\) and \(\operatorname{Re} s > 1\), let \(\zeta(s, x) := \sum_{n \ge 0} (n + x)^{-s}\).  Euler--Maclaurin summation gives
\[
  \zeta(s, x) = \frac{x^{1-s}}{s - 1} + \frac{x^{-s}}{2} - s \int_0^\infty \frac{\{t\} - \frac12}{(t + x)^{s+1}} \, d t ,
\]
which continues \(\zeta(s, x)\) to \(\operatorname{Re} s > -1\) with a single pole at \(s = 1\).  Write \(\zeta'(s, x)\) for \(\partial_s \zeta(s, x)\).

\begin{problem}\label{pr:hz0}
\(\zeta(0, x) = \frac12 - x\), and \(\zeta(s, \frac12) = (2^s - 1)\zeta(s)\).
\end{problem}

\begin{problem}\label{pr:lerch-diff}
\(\partial_x^2\, \zeta'(0, x) = \zeta(2, x) = (\log\Gamma)''(x)\), and \(\zeta'(0, x+1) - \zeta'(0, x) = \log x\).
\end{problem}

Problems~\ref{pr:hz0} and~\ref{pr:lerch-diff}, with \(\zeta(0) = -\frac12\) and \(\Gamma(\frac12) = \sqrt\pi\), prove the following.

\begin{theorem}[Lerch]\label{thm:lerch}
\(\zeta'(0, x) = \log\Gamma(x) - \frac12\log 2\pi\) for \(x > 0\).
\end{theorem}

\begin{problem}\label{pr:L0}
Let \(\chi\) be an odd character modulo \(N\).  Then \(L(0, \chi) = -\frac1N \sum_{a=1}^{N} a\,\chi(a)\), and
\[
  L'(0, \chi) = \sum_{a=1}^{N-1} \chi(a) \log\Gamma(a/N) - L(0, \chi) \log N .
\]
\end{problem}

Let \(\chi\) be a primitive odd character modulo \(N\), and let
\[
  \Lambda(s, \chi) := (N/\pi)^{(s+1)/2}\, \Gamma\bigl(\tfrac{s+1}{2}\bigr) L(s, \chi) .
\]  Then \(\Lambda(s, \chi) = W(\chi)\Lambda(1 - s, \bar\chi)\) with \(\lvert W(\chi)\rvert = 1\) \cite[Ch.~9]{Dav}; moreover \(W(\chi) = \tau(\chi)/(\ii\sqrt N)\), which is \(1\) if \(\chi\) is real, because then the Gauss sum is \(\tau(\chi) = \ii\sqrt N\) \cite[Ch.~6]{Dav}.

\begin{problem}\label{pr:FE}
\(\dfrac{L'(1, \chi)}{L(1, \chi)} + \dfrac{L'(0, \bar\chi)}{L(0, \bar\chi)} = \gamma + \log\dfrac{2\pi}{N}\).
\end{problem}

\begin{problem}\label{pr:LL4}
For the character modulo \(4\), Problems~\ref{pr:L0} and~\ref{pr:FE} give Theorem~\ref{ec-thm:LL}:
\[
  \frac{L'(1, \chi)}{L(1, \chi)} = \gamma + 2\log 2 + 3\log\pi - 4\log\Gamma(\tfrac14) = 0.24560\ldots
\]
\end{problem}

\begin{problem}\label{pr:eta-i}
\(\eta(\ii) = \Gamma(\frac14)/(2\pi^{3/4})\).
\end{problem}

Chapter~\ref{ec-chap} found this value from \(\Delta(\varpi\Z[\ii]) = 64\) (Problem~\ref{ec-pr:eta-i}).  Here it comes from computing one constant, \(L'(1,\chi)/L(1,\chi)\), in two ways: through \(s = 1\) by Theorem~\ref{thm:klf}, and through \(s = 0\) by Theorem~\ref{thm:lerch}.  The argument used only that \(Z(s)\) is both an Eisenstein series and \(4\zeta(s)L(s,\chi)\); the next section runs it for every imaginary quadratic field.

\section{The general odd character}\label{sec:general}

Let \(K\) be an imaginary quadratic field of discriminant \(-d\), with class number \(h\) and \(w\) roots of unity, and let \(\chi\) be its character, a primitive real odd character modulo \(d\).  The class number formula is \(L(1, \chi) = 2\pi h/(w\sqrt d)\).  For each class of positive definite forms of discriminant \(-d\), choose the reduced representative \(Q = [a, b, c]\), so that \(Q(m,n) = am^2 + bmn + cn^2\), \(b^2 - 4ac = -d\), \(\lvert b\rvert \le a \le c\), and \(b \ge 0\) if \(\lvert b\rvert = a\) or \(a = c\).  Let \(\tau_Q := (-b + \ii\sqrt d)/2a\) and \(y_Q := \operatorname{Im}\tau_Q\).  The forms \(Q\) together represent each \(n \ge 1\) exactly \(w \sum_{e \mid n} \chi(e)\) times \cite[Ch.~6]{Dav}, so
\[
  \sum_Q \sum_{(m,n) \ne (0,0)} Q(m,n)^{-s} = w\, \zeta(s) L(s, \chi)
  \qquad (\operatorname{Re} s > 1) .
\]
The sum over \(Q\) has \(h\) terms.  When \(h > 1\), can one class be taken alone?

\begin{problem}\label{pr:d20-reps}
For \(d = 20\) the reduced forms are \(Q_1 = [1, 0, 5]\) and \(Q_2 = [2, 2, 3]\).  The form \(Q_1\) represents \(1\) twice and \(3\) not at all, \(Q_2\) represents \(1\) not at all and \(3\) four times, and \(w\sum_{e \mid n}\chi(e)\) is \(2\) for \(n = 1\) and \(4\) for \(n = 3\).  Hence neither \(\sum'_{(m,n)} Q_i(m,n)^{-s}\) is a constant multiple of \(\zeta(s)L(s, \chi)\).
\end{problem}

So a formula through \(L(s, \chi)\) can see only the sum over classes.  At \(s = 1\) this sum is a sum of logarithms, and what one can hope to evaluate is the product over \(Q\) of the quantities \(\sqrt{y_Q}\,\lvert\eta(\tau_Q)\rvert^2\), not each one.

\begin{problem}\label{pr:L0-gen}
\(L(0, \chi) = 2h/w\).  For \(d = 3, 4, 7, 23\) this is \(\frac13, \frac12, 1, 3\).
\end{problem}

\begin{problem}\label{pr:lerch-gen}
\(\displaystyle \frac{L'(1, \chi)}{L(1, \chi)} = \gamma + \log 2\pi - \frac{w}{2h} \sum_{a=1}^{d-1} \chi(a) \log\Gamma\Bigl(\frac ad\Bigr)\).  For \(d = 3, 7, 23\) it is \(0.36828\ldots\), \(0.01563\ldots\) and \(-0.42212\ldots\).
\end{problem}

\begin{problem}\label{pr:form-E}
\(Q(m, n) = \dfrac{\sqrt d}{2 y_Q}\lvert m - n\tau_Q\rvert^2\), and so
\[
  \sum_Q E(\tau_Q, s) = w \Bigl(\frac{\sqrt d}{2}\Bigr)^{s} \zeta(s) L(s, \chi) .
\]
\end{problem}

\begin{problem}\label{pr:LL-eta-gen}
\(\displaystyle \frac{L'(1, \chi)}{L(1, \chi)} = \gamma - \log 2 - \frac12\log d - \frac2h \sum_Q \log\bigl(\sqrt{y_Q}\, \lvert\eta(\tau_Q)\rvert^2\bigr)\).
\end{problem}

Problems~\ref{pr:lerch-gen} and~\ref{pr:LL-eta-gen} prove the formula of Chowla and Selberg \cite{CS}.

\begin{theorem}[Chowla--Selberg]\label{thm:cs}
\[
  \prod_Q \sqrt{y_Q}\, \lvert\eta(\tau_Q)\rvert^2 = (4\pi\sqrt d)^{-h/2} \prod_{a=1}^{d-1} \Gamma\Bigl(\frac ad\Bigr)^{w\chi(a)/4} .
\]
\end{theorem}

\begin{problem}\label{pr:d20-mean}
For \(d = 20\), \(y_{Q_1} = \sqrt5\) and \(y_{Q_2} = \sqrt5/2\).  The geometric mean of \(\sqrt{y_Q}\,\lvert\eta(\tau_Q)\rvert^2\) over the two classes is
\[
  (8\pi\sqrt5)^{-1/2} \Bigl(\frac{\Gamma(\frac1{20})\Gamma(\frac3{20})\Gamma(\frac7{20})\Gamma(\frac9{20})}{\Gamma(\frac{11}{20})\Gamma(\frac{13}{20})\Gamma(\frac{17}{20})\Gamma(\frac{19}{20})}\Bigr)^{1/4} = 0.52301\ldots ,
\]
while the two values themselves are \(0.46373\ldots\) for \(Q_1\) and \(0.58987\ldots\) for \(Q_2\).
\end{problem}

\begin{problem}\label{pr:eta-rho}
Let \(\rho := e^{2\pi\ii/3}\).  Then \(\lvert\eta(\rho)\rvert^2 = 3^{1/4}\, \Gamma(\frac13)^3/4\pi^2\).
\end{problem}

\section{Periods}\label{sec:periods}

Why should \(\prod_a \Gamma(a/d)^{\chi(a)}\) be a period?  Because values of \(\Gamma\) at rational numbers come, up to algebraic numbers and powers of \(\pi\), from integrals of algebraic differentials on Fermat curves.

Let \(N \ge 3\), and let \(F_N\) be the smooth projective curve with affine equation \(u^N + v^N = 1\).  For integers \(a, b \ge 1\) with \(a + b \le N - 1\), the differential \(\omega_{a,b} := u^{a-1}v^{b-N}\, d u\) is holomorphic on \(F_N\), as one checks at the points with \(v = 0\) and at infinity.  Let \(\kappa\) be the path \(u \mapsto (u, (1 - u^N)^{1/N})\), \(0 \le u \le 1\), from \((0,1)\) to \((1,0)\).  Let \(\zeta_N := e^{2\pi\ii/N}\), and let \(A(u,v) := (\zeta_N u, v)\) and \(B(u,v) := (u, \zeta_N v)\).

\begin{problem}\label{pr:beta-path}
\(\int_\kappa \omega_{a,b} = \frac1N \Beta(\frac aN, \frac bN)\).
\end{problem}

\begin{problem}\label{pr:cycle}
The chain \((1 - A)(1 - B)\kappa\) is a closed cycle on \(F_N\), and the integral of \(\omega_{a,b}\) over it is \(\frac1N (1 - \zeta_N^a)(1 - \zeta_N^b)\, \Beta(\frac aN, \frac bN)\).
\end{problem}

So each \(\Beta(\frac aN, \frac bN)\) with \(a, b \ge 1\) and \(a + b \le N - 1\) is, up to an algebraic factor, a period of a holomorphic differential on \(F_N\) defined over \(\Q\).

\begin{problem}\label{pr:gamma-N}
For integers \(a \ge 1\), \(\Gamma(\frac aN)^N = \Gamma(a) \prod_{j=1}^{N-1} \Beta(\frac aN, \frac{ja}N)\).
\end{problem}

Since \(\Beta(x, y+1) = \Beta(x, y)\, y/(x+y)\), every \(\Beta(\frac aN, \frac bN)\) with \(a, b \ge 1\) is a rational multiple of one with \(1 \le a, b \le N\), which is \(N\) times an integral of the algebraic differential \(\omega_{a,b}\) along \(\kappa\) (Problem~\ref{pr:beta-path}).  So by Problem~\ref{pr:gamma-N}, \(\Gamma(\frac aN)^N\) is, up to a rational factor, a product of such integrals.

\begin{problem}\label{pr:quartic}
The map \((u, v) \mapsto (u, v^2)\) from \(F_4\) to the curve \(y^2 = 1 - t^4\) pulls \(d t / y\) back to \(\omega_{1,2}\), and \(\Beta(\frac14, \frac12) = 2\varpi\).
\end{problem}

For \(K\) as in Section~\ref{sec:general}, let
\[
  \Omega_K := \sqrt\pi \prod_{a=1}^{d-1} \Gamma\Bigl(\frac ad\Bigr)^{w\chi(a)/4h} .
\]
For \(K = \Q(\ii)\), \(\Omega_K = 2\varpi\) (Problem~\ref{ec-pr:Omega}, where \(\Omega\) denotes \(\Omega_K^2/2\pi = 2\varpi^2/\pi\)), which by Problems~\ref{pr:cycle} and~\ref{pr:quartic} is a period of \(d t / y\) on \(y^2 = 1 - t^4\) up to a factor in \(\bar\Q^\times\).

\begin{problem}\label{pr:Omega3}
For \(K = \Q(\sqrt{-3})\),
\[
  \Omega_K = \sqrt\pi\, \bigl(\Gamma(\tfrac13)/\Gamma(\tfrac23)\bigr)^{3/2} = 3^{1/4}2^{-1/2}\, \Beta(\tfrac13, \tfrac13) .
\]
\end{problem}

Here \(F_3\) is itself a curve of genus one, \(\omega_{1,1}\) is its holomorphic differential, and \(A\) multiplies \(\omega_{1,1}\) by \(\zeta_3\): the curve has complex multiplication by \(\Z[\rho]\).  This is the general pattern.  Chowla and Selberg deduced the following from Theorem~\ref{thm:cs} \cite{CS}; we give it in the formulation of Gross \cite{Gr}.  Every nonzero period of a nonzero holomorphic differential defined over \(\bar\Q\) on an elliptic curve over \(\bar\Q\) with complex multiplication by \(\Oc_K\) lies in \(\bar\Q^\times\, \Omega_K\).  Gross proved this period relation geometrically, by relating such curves to Jacobians of Fermat curves, whose periods are Beta values by Problem~\ref{pr:cycle} \cite{Gr}.

\begin{problem}\label{pr:cm-period}
\(e^{\gamma - L'(1,\chi)/L(1,\chi)} = \Omega_K^2/2\pi^2\).  For \(K = \Q(\ii)\) this is \(2\varpi^2/\pi^2 = \Gamma(\frac14)^4/4\pi^3 = 1.39320\ldots\)
\end{problem}

Thus \(e^{\gamma - L'(1,\chi)/L(1,\chi)} = -2\,(\Omega_K/2\pi\ii)^2\): up to a factor in \(\bar\Q^\times\), it is the square of the ratio of a period of an elliptic curve with complex multiplication by \(\Oc_K\) to \(2\pi\ii\), the period of \(d z/z\) on the multiplicative group.  In Chapter~\ref{cp-chap}, \emph{Covolumes and periods}, the two constants enter in this form for \(K = \Q(\ii)\), through the Faltings height of such a curve (Problem~\ref{cp-pr:hF-CS}).

\section*{Hints}

\subsection*{\S\ref{sec:klf}}
\begin{description}
\item[\ref{pr:E-inv}] For \(g = \bigl(\begin{smallmatrix} a & b \\ c & d \end{smallmatrix}\bigr)\), write \(m + n\, g\tau\) as \((m' + n'\tau)/(c\tau + d)\), and use \(\operatorname{Im}(g\tau) = y/\lvert c\tau + d\rvert^2\).
\item[\ref{pr:poisson}] Apply Poisson summation in \(m\).  To find the Fourier transform of \((t^2 + y^2)^{-s}\), write it as an integral of Gaussians in \(t\) by means of \(\Gamma(s)\).
\item[\ref{pr:fourier}] Separate \(n = 0\).  For \(n \ne 0\) apply Problem~\ref{pr:poisson} with \(nx\) and \(\lvert n\rvert y\), and collect the terms with \(kn = N\).  For the continuation, \(K_\nu(z)\) decays exponentially as \(z \to \infty\), uniformly for \(\nu\) in a compact set.
\item[\ref{pr:K-half}] Substitute \(t = e^{2v}\) and then \(w = \sinh v\), and reduce to a Gaussian integral.
\item[\ref{pr:q-terms}] \(\sum_{n \ge 1} \sigma_{-1}(n)\, q^n = -\sum_{m \ge 1}\log(1 - q^m)\).
\item[\ref{pr:const}] \(\zeta(2s - 1) = \frac{1}{2(s-1)} + \gamma + O(s-1)\), and \(\psi(\frac12) = -\gamma - 2\log 2\) by the duplication formula.
\item[\ref{pr:eta-mod}] Theorem~\ref{thm:klf} and Problem~\ref{pr:E-inv}.
\item[\ref{pr:Zi}] \(Z(s) = E(\ii, s)\) and \(\eta(\ii) > 0\).  Then compare with the Laurent expansion of \(4\zeta(s)L(s,\chi)\), using \(L(1, \chi) = \pi/4\).
\end{description}

\subsection*{\S\ref{sec:lerch}}
\begin{description}
\item[\ref{pr:hz0}] Set \(s = 0\) in the continuation formula.  For the second statement, split \(\zeta(s)\) into even and odd \(n\).
\item[\ref{pr:lerch-diff}] \(\partial_x \zeta(s, x) = -s\, \zeta(s+1, x)\).  Differentiate twice in \(x\), then in \(s\) at \(s = 0\).  For the second statement, \(\zeta(s, x+1) = \zeta(s, x) - x^{-s}\).
\item[\ref{pr:L0}] \(L(s, \chi) = N^{-s}\sum_{a=1}^{N} \chi(a)\, \zeta(s, a/N)\), and \(\sum_a \chi(a) = 0\).
\item[\ref{pr:FE}] Take the logarithmic derivative of the functional equation at \(s = 1\), and use \(\psi(\frac12)\) as in Problem~\ref{pr:const}.  Here \(L(0, \bar\chi) \ne 0\), since \(\Lambda(0, \bar\chi)\) is a nonzero multiple of \(\Lambda(1, \chi)\).
\item[\ref{pr:LL4}] \(\Gamma(\frac14)\Gamma(\frac34) = \pi\sqrt2\).
\end{description}

\subsection*{\S\ref{sec:general}}
\begin{description}
\item[\ref{pr:L0-gen}] Compare \(\Lambda(0, \chi)\) with \(\Lambda(1, \chi)\), and use the class number formula.
\item[\ref{pr:lerch-gen}] Problems~\ref{pr:L0}, \ref{pr:FE} and~\ref{pr:L0-gen}.
\item[\ref{pr:form-E}] Expand \(a\lvert m - n\tau_Q\rvert^2\).  Then replace \(n\) by \(-n\).
\item[\ref{pr:LL-eta-gen}] Expand both sides of Problem~\ref{pr:form-E} at \(s = 1\) by Theorem~\ref{thm:klf}.  The residues agree by the class number formula; compare the constant terms.
\item[\ref{pr:d20-reps}] List the small values of each form.  For the right side, \(\chi(3) = (\frac{-20}{3})\).
\item[\ref{pr:d20-mean}] Theorem~\ref{thm:cs} with \(h = w = 2\); list the \(a < 20\) with \(\chi(a) = 1\).  For the numerical values, sum the product for \(\eta\) to a few terms.
\item[\ref{pr:eta-rho}] Theorem~\ref{thm:cs} with \(d = 3\), and the reflection formula.
\end{description}

\subsection*{\S\ref{sec:periods}}
\begin{description}
\item[\ref{pr:beta-path}] Substitute \(t = u^N\) in \(\Beta(\frac aN, \frac bN) = \int_0^1 t^{a/N - 1}(1-t)^{b/N - 1}\, d t\).
\item[\ref{pr:cycle}] \(A^*\omega_{a,b} = \zeta_N^a\, \omega_{a,b}\) and \(B^*\omega_{a,b} = \zeta_N^b\, \omega_{a,b}\).  For the boundary, \(A\) fixes the points with \(u = 0\) and \(B\) those with \(v = 0\).
\item[\ref{pr:gamma-N}] \(\Beta(x, y) = \Gamma(x)\Gamma(y)/\Gamma(x+y)\); the product telescopes.
\item[\ref{pr:quartic}] Problem~\ref{pr:beta-path} and the definition of \(\varpi\) in Chapter~\ref{ec-chap}.
\item[\ref{pr:Omega3}] \(h = 1\), \(w = 6\), and the reflection formula as in Problem~\ref{pr:eta-rho}.
\item[\ref{pr:cm-period}] Problem~\ref{pr:lerch-gen}.
\end{description}

\begin{thebibliography}{CS}
\bibitem[Ap]{Ap} T.~M.~Apostol, \emph{Introduction to Analytic Number Theory}, Undergrad.\ Texts Math., Springer, 1976.
\bibitem[CS]{CS} A.~Selberg and S.~Chowla, \emph{On Epstein's zeta-function}, J.\ reine angew.\ Math.\ \textbf{227} (1967), 86--110.
\bibitem[Dav]{Dav} H.~Davenport, \emph{Multiplicative Number Theory}, 3rd ed., Grad.\ Texts in Math.\ 74, Springer, 2000.
\bibitem[Gr]{Gr} B.~H.~Gross, \emph{On the periods of abelian integrals and a formula of Chowla and Selberg}, with an appendix by D.~E.~Rohrlich, Invent.\ Math.\ \textbf{45} (1978), 193--211.
\bibitem[La]{La} S.~Lang, \emph{Elliptic Functions}, 2nd ed., Grad.\ Texts in Math.\ 112, Springer, 1987.
\bibitem[We]{We} A.~Weil, \emph{Elliptic Functions according to Eisenstein and Kronecker}, Ergeb.\ Math.\ Grenzgeb.\ 88, Springer, 1976.
\end{thebibliography}

\end{document}
